\chapter{Phlebitis}

\section*{Definition}
Phlebitis is inflammation of a vein, classified as superficial thrombophlebitis (involving superficial veins with associated thrombosis) or phlebitis without thrombosis. Superficial thrombophlebitis presents as a tender, erythematous, palpable cord along the course of a superficial vein.

\section*{What Happens in Your Body}
Endothelial inflammation leads to platelet aggregation, fibrin deposition, and thrombus formation within the vein lumen. The inflammatory response causes localized pain, erythema, warmth, and induration; extension into deep veins or embolization occurs in a minority of superficial cases.

\section*{Causes}
\begin{itemize}
  \item \textbf{Iatrogenic:} Peripheral IV catheter (most common cause), intravenous medications (chemotherapy, potassium, vancomycin), extravasation of irritant solutions
  \item \textbf{Mechanical:} Venous stasis, varicose veins, trauma, repetitive irritation
  \item \textbf{Infectious:} Septic phlebitis from IV catheter colonization or local infection (Staphylococcus aureus, Gram-negative bacilli), Lemierre syndrome (thrombophlebitis of internal jugular vein)
  \item \textbf{Hypercoagulable:} Malignancy (Trousseau syndrome), pregnancy, oral contraceptives, factor V Leiden, antiphospholipid syndrome
  \item \textbf{Other:} Beh\c{c}et disease, Buerger disease (thromboangiitis obliterans), migratory superficial thrombophlebitis
\end{itemize}

\section*{When to See a Doctor}
See your doctor if:
\begin{itemize}
  \item Extensive or rapidly spreading erythema along a vein
  \item Fever, chills, purulent drainage from IV site (suggests septic phlebitis)
  \item Vein involvement above the knee (greater saphenous) with risk of deep vein extension
  \item Known hypercoagulable state or malignancy
  \item Recurrent or migratory phlebitis without identifiable cause
\end{itemize}

\section*{Related Symptoms}
\begin{itemize}
  \item Pain, tenderness, erythema, warmth, palpable cord, swelling along the course of a superficial vein, low-grade fever
\end{itemize}
\textbf{Important:} While typically self-limited, superficial thrombophlebitis of the great saphenous vein above the knee carries a 10--15\% risk of extension into the deep femoral vein.

\section*{Self-Care / Home Management}
\begin{itemize}
  \item Warm compresses or heat packs applied to the affected area for 15--20 minutes, 3--4 times daily
  \item NSAIDs (ibuprofen 400--600 mg or diclofenac topical gel) for pain and inflammation
  \item Elevation of the affected limb to promote venous drainage and reduce edema
  \item Avoid prolonged standing or tight clothing over the affected area
  \item Continue physical activity; bed rest is not required for superficial phlebitis
\end{itemize}

\section*{How Your Doctor Will Evaluate You}
\begin{itemize}
  \item Clinical diagnosis based on tender, erythematous cord along a superficial vein
  \item Compression ultrasound to differentiate superficial thrombophlebitis from DVT and assess extent of involvement
  \item Doppler ultrasound to evaluate for deep vein extension, especially in above-knee greater saphenous involvement
  \item Culture of purulent discharge or catheter tip if septic phlebitis suspected
  \item Hypercoagulability workup for recurrent or unprovoked phlebitis (factor V Leiden, prothrombin mutation, antiphospholipid antibodies)
\end{itemize}

\section*{Treatment Options}
\begin{itemize}
  \item Remove IV catheter if phlebitis is catheter-related; apply warm compresses
  \item NSAIDs (oral or topical) as first-line therapy for pain and inflammation
  \item Compression stockings for lower extremity involvement with underlying venous insufficiency
  \item Anticoagulation (LMWH or fondaparinux) for 45 days if superficial thrombophlebitis extends more than 5 cm and involves the great saphenous vein above the knee
  \item Antibiotics (vancomycin or cefazolin) for septic phlebitis with IV catheter removal and possible surgical excision
\end{itemize}

\section*{References}
\begin{thebibliography}{99}
\bibitem{ref1} Nasr H, Scriven JM. "Superficial thrombophlebitis: a review of the literature." \textit{European Journal of Vascular and Endovascular Surgery}, 2015;50(6):740--746.
\bibitem{ref2} Decousus H, Qu\'er\'e I, Presles E, et al. "Superficial venous thrombosis and venous thromboembolism: a large, prospective epidemiologic study." \textit{Annals of Internal Medicine}, 2010;152(4):218--224.
\bibitem{ref3} Di Nisio M, Wichers IM, Middeldorp S. "Treatment of superficial thrombophlebitis." \textit{Cochrane Database of Systematic Reviews}, 2018;(3):CD004982.
\bibitem{ref4} Gallerani M, Manfredini R, Portaluppi F, et al. "Risk factors for superficial vein thrombosis." \textit{International Angiology}, 2005;24(2):164--169.
\bibitem{ref5} Kalodiki E, Stvrtinova V, Allegra C, et al. "Superficial vein thrombosis: a consensus document." \textit{International Angiology}, 2012;31(3):203--216.
\bibitem{ref6} Beyer-Westendorf J, Bouill{\`e}gue F, Gremillet E, et al. "Diagnosis and management of superficial vein thrombosis." \textit{H{\"a}mostaseologie}, 2017;37(1):37--45.

\end{thebibliography}

\clearpage
